\documentclass[12pt,a4paper]{report}
\usepackage[utf8]{inputenc}
\usepackage[spanish]{babel}
\usepackage[margin=2.5cm]{geometry}
\usepackage{graphicx}
\usepackage{xcolor}
\usepackage{hyperref}
\usepackage{enumitem}
\usepackage{listings}
\usepackage{tabularx}
\usepackage{booktabs}

% Configuración de listings para mostrar correctamente código PHP
\setcounter{tocdepth}{2}
\lstset{
  language=PHP,
  basicstyle=\ttfamily\small,
  numbers=left,
  numberstyle=\tiny\color{gray},
  stepnumber=1,
  numbersep=5pt,
  breaklines=true,
  breakatwhitespace=true,
  postbreak=\mbox{\textcolor{magentaDark}{$\hookrightarrow$}\space},
  columns=fullflexible,
  keepspaces=true,
  showstringspaces=false,
  tabsize=4
}

% Colores magenta
\definecolor{magentaMain}{HTML}{FF007F}
\definecolor{magentaDark}{HTML}{C80064}
\hypersetup{colorlinks=true,linkcolor=magentaDark,urlcolor=magentaMain}

\title{Documentación Técnica: TAD E-commerce Magenta}
\author{Pablo Saez Morales \\ Oleksandr Baranets \\ Nicolas Porra Collado}
\date{10 de Mayo de 2026}

\begin{document}

\maketitle
\begin{center}
  \vspace{0.5cm}
  Repositorio: https://github.com/pablosaez21/TAD-Ecommerce\\
  Grupo Magenta: Pablo Saez Morales, Oleksandr Baranets, Nicolas Porra Collado
\end{center}

\tableofcontents
\cleardoublepage
% Nota: compile el fichero dos veces si el índice no aparece actualizado.
% Ejemplo: lualatex DOCUMENTACION.tex && lualatex DOCUMENTACION.tex

\chapter{Resumen ejecutivo}

TAD E-commerce Magenta es una aplicación de comercio electrónico desarrollada con \textbf{Laravel 9} que soporta catálogo de productos, carrito en sesión, proceso de checkout con transacciones ACID, gestión de direcciones, panel administrativo y sistema multiidioma (es/en). Esta documentación recoge diseño, decisiones técnicas, estructura, instrucciones de instalación y uso, pruebas y conclusiones.

\chapter{Objetivos y alcance}

\section{Objetivos}
\begin{itemize}
  \item Implementar un e-commerce funcional con autenticación segura (Fortify).
  \item Asegurar coherencia de datos en el checkout mediante transacciones y bloqueo de stock.
  \item Proveer un panel administrativo para gestión de catálogo y pedidos.
  \item Mantener buenas prácticas de seguridad y rendimiento.
\end{itemize}

\section{Alcance}

El alcance cubre las funcionalidades implementadas en el repositorio actual: modelos, migraciones, controladores, vistas Blade, rutas, internacionalización y envíos de correo (Mailables). No incluye integraciones de terceros en producción (pasarelas de pago reales) aunque la arquitectura permite su incorporación.

\chapter{Arquitectura y diseño}

\section{Visión general}

La aplicación sigue el patrón MVC de Laravel: modelos Eloquent para la lógica de datos, controladores para orquestación, vistas Blade para presentación. Los assets se gestionan con Vite y npm.

\section{Modelos principales}

Resumen de modelos y relaciones (implementadas en \texttt{app/Models}):

\begin{itemize}
  \item \textbf{User}: relaciones \texttt{hasMany(addresses, orders)}, \texttt{belongsToMany(products) via favorites}.
  \item \textbf{Product}: relaciones \texttt{belongsToMany(categories)}, \texttt{hasMany(orderLines)}, campos: \texttt{name, description, price, stock, active}.
  \item \textbf{Category}: \texttt{belongsToMany(products)}.
  \item \textbf{Order}: \texttt{belongsTo(user,address)}, \texttt{hasMany(orderLines)}, \texttt{hasOne(payment)}.
  \item \textbf{OrderLine}: \texttt{belongsTo(order,product)}, guarda \texttt{quantity, unit_price, subtotal}.
  \item \textbf{Address}: \texttt{belongsTo(user)} con campos de dirección y \texttt{is_default}.
  \item \textbf{Payment}: \texttt{belongsTo(order)} registra el método y estado.
\end{itemize}

\subsection{Detalles y extractos de código}

Los siguientes extractos provienen de los ficheros reales en \texttt{app/Models} y ayudan a entender decisiones concretas.

\subsubsection{Producto (\texttt{app/Models/Product.php})}
El modelo define los campos rellenables y las relaciones principales:
\begin{lstlisting}[language=PHP]
protected $fillable = ['name','description','price','image','stock','active'];

public function categories(): BelongsToMany { return $this->belongsToMany(Category::class); }
public function orderLines(): HasMany { return $this->hasMany(OrderLine::class); }
public function favoredByUsers(): BelongsToMany { return $this->belongsToMany(User::class, 'favorites'); }
\end{lstlisting}

Notas:
\begin{itemize}
  \item El campo \texttt{price} se castea a decimal con dos decimales para evitar errores de formato.
  \item La relación \texttt{favoredByUsers} usa la tabla pivote \texttt{favorites}.
\end{itemize}

\subsubsection{Orden (\texttt{app/Models/Order.php})}
El modelo de orden relaciona líneas, usuario, dirección y pago:
\begin{lstlisting}[language=PHP]
protected $fillable = ['user_id','address_id','status','total'];

public function lines(): HasMany { return $this->hasMany(OrderLine::class); }
public function payment(): HasOne { return $this->hasOne(Payment::class); }
\end{lstlisting}

Observaciones:
\begin{itemize}
  \item El total se almacena con precisión decimal (\texttt{decimal:2}).
  \item Se mantiene copia de precio en \texttt{order\_lines} para auditoría histórica.
\end{itemize}

\section{Esquema de base de datos}

Se aplica normalización hasta 3NF; las tablas pivot eliminan redundancia en relaciones N:M (ej. \texttt{favorites}, \texttt{category\_product}, \texttt{discount\_product}). Las decisiones principales:

\begin{itemize}
  \item Precios se copian en \texttt{order\_lines} para auditoría histórica.
  \item Las claves foráneas usan \texttt{onDelete('cascade')} cuando procede.
  \item Índices en columnas de búsqueda (\texttt{created\_at}, \texttt{status}).
\end{itemize}

\chapter{Rutas y controladores}

\section{Rutas principales}

Rutas definidas en \texttt{routes/web.php} (resumen):

\begin{itemize}
  \item PÚBLICAS: \texttt{GET /}, \texttt{GET /products}, \texttt{GET /products/{product}}, \texttt{GET /categories/{slug}}.
  \item AUTENTICADAS: \texttt{GET/POST /cart}, \texttt{POST /cart/checkout}, perfil y direcciones, favoritos.
  \item ADMIN: rutas protegidas con middleware \texttt{admin} para \texttt{/admin} y CRUD de productos y categorías.
\end{itemize}

\section{Controladores clave}

\begin{itemize}
  \item \textbf{ProductController}: list, show, create, store, edit, update, destroy (admin + público).
  \item \textbf{CartController}: gestionar carrito en sesión, validar stock, checkout (transacción).
  \item \textbf{ProfileController}: gestionar direcciones, idioma, ver pedidos.
  \item \textbf{FavoriteController}: toggle favoritos.
  \item \textbf{DashboardController}: estadísticas administrativas.
\end{itemize}

\subsection{CartController: métodos clave}

Se muestran extractos reales de \texttt{app/Http/Controllers/CartController.php} y una explicación del comportamiento:

\subsubsection{Añadir al carrito (\texttt{add})}
El método valida disponibilidad del producto y mantiene el carrito en sesión. Fragmento:
\begin{lstlisting}[language=PHP]
if (! $product->active || $product->stock <= 0) {
    return back()->with('error', 'Este producto no está disponible.');
}

$cart = session('cart', []);
$cart[$product->id] = [... 'quantity' => min($quantity, $product->stock) ...];
session(['cart' => $cart]);
\end{lstlisting}

Notas:
\begin{itemize}
  \item El carrito se guarda como array asociativo en la sesión: clave = id del producto.
  \item La cantidad se clampa por el stock disponible para evitar solicitudes inválidas.
\end{itemize}

\subsubsection{Checkout (\texttt{checkout})}
El checkout es el núcleo del flujo de negocio: usa validación, transacción y bloqueo pessimistico para stock.
\begin{lstlisting}[language=PHP]
DB::transaction(function () use ($validated, $user, $address) {
    $order = Order::create([...]);
    foreach ($validated['items'] as $item) {
        $product = Product::whereKey($item['product_id'])
            ->where('active', true)
            ->lockForUpdate()
            ->firstOrFail();

        if ($product->stock < $item['quantity']) {
            abort(422, "Stock insuficiente para {$product->name}");
        }

        $order->lines()->create([ ... ]);
        $product->decrement('stock', $item['quantity']);
    }
    $order->payment()->create([...]);
    return $order->fresh(['lines.product', 'payment', 'address']);
});
\end{lstlisting}

Explicación:
\begin{itemize}
  \item Se usa \texttt{lockForUpdate()} para bloquear la fila del producto durante la transacción.
  \item Si falta stock, se aborta con código 422, provocando rollback.
  \item La creación del pago se realiza dentro de la misma transacción para mantener consistencia.
\end{itemize}

\chapter{Proceso de compra (checkout)}

\section{Flujo normal}

\begin{enumerate}
  \item Usuario autenticado inicia checkout desde carrito en sesión.
  \item Sistema valida stock de cada producto con bloqueo (\texttt{lockForUpdate}).
  \item Se crea la orden, order\_lines y registro de payment dentro de una \texttt{DB::transaction}.
  \item Se decrementa stock y se actualiza total.
  \item Se envía correo de confirmación (Mailable) y se limpia el carrito.
\end{enumerate}

\section{Manejo de errores y rollback}

Cualquier fallo en la transacción provoca rollback. Casos alternativos: carrito vacío, stock insuficiente, fallo en envío de correo (reintentos asincrónicos / logging).

\chapter{Decisiones técnicas y de diseño}

\section{Autenticación}

Se eligió \textbf{Laravel Fortify} para autenticación y gestión de credenciales; permite 2FA y acciones personalizadas (CreateNewUser, UpdateUserPassword, etc.).

\section{Carrito en sesión}

Ventaja: simple, rápida y suficiente para la EPD. Alternativa (persistente en BD) considerada para futuras versiones.

\section{Transacciones y concurrencia}

Uso de transacciones ACID y \texttt{lockForUpdate} para evitar sobreventa de stock en escenarios concurrentes.

\section{Internacionalización}

Se utilizan archivos PHP en \texttt{resources/lang/es} y \texttt{resources/lang/en}. La preferencia de idioma se guarda en el campo \texttt{language} del usuario.

\chapter{Estructura de ficheros relevante}

\begin{itemize}
  \item \texttt{app/Models/}: modelos Eloquent.
  \item \texttt{app/Http/Controllers/}: controladores.
  \item \texttt{resources/views/}: plantillas Blade (productos, carrito, checkout, perfil, admin).
  \item \texttt{routes/web.php}: rutas web.
  \item \texttt{database/migrations/}: migraciones y seeders.
  \item \texttt{app/Mail/}: clases Mailable (OrderConfirmationMail).
\end{itemize}

\chapter{Instalación y puesta en marcha}

\section{Requisitos}

\begin{itemize}
  \item PHP 8.1+
  \item Composer
  \item Node.js 16+ y npm
  \item MySQL o PostgreSQL
\end{itemize}

\section{Pasos rápidos}

\begin{verbatim}
# clona el repo (local ya presente)
cd /home/miller/Desktop/tad/TAD-Ecommerce
composer install
cp .env.example .env
php artisan key:generate
# configurar .env con credenciales DB y MAIL
php artisan migrate --seed
npm install
npm run dev
php artisan serve

over: http://localhost:8000
\end{verbatim}

\chapter{Configuración de correo}

Se recomienda Mailtrap para desarrollo o configuraciones SMTP en .env (MAIL_HOST, MAIL_PORT, MAIL_USERNAME, MAIL_PASSWORD).

\chapter{Pruebas}

\section{Pruebas unitarias y funcionales}

Tests incluidos en \texttt{tests/Feature} y \texttt{tests/Unit} cubren: modelos, checkout, creación de órdenes y validaciones básicas. Ejecutar:

\begin{verbatim}
php artisan test
\end{verbatim}

\chapter{Seguridad}

\begin{itemize}
  \item Contraseñas cifradas (bcrypt).
  \item Protección CSRF en formularios (Blade @csrf).
  \item Validación en servidores mediante FormRequest.
  \item Middleware de autorización (admin) para rutas sensibles.
\end{itemize}

\chapter{Eliminación de contenido irrelevante}

Se han eliminado todas las referencias a la documentación y recursos del proyecto previo ("Sabores de Barrio"). No quedan inclusiones de imágenes externas que pertenecieran a ese proyecto. Si existen ficheros de imagen en el repositorio que no pertenecen a TAD E-commerce y deseas que los borre del disco, indícalo y procederé con cuidado.

\chapter{Decisiones de mejoras futuras}

\begin{itemize}
  \item Integrar pasarela de pago real (Stripe/PayPal).
  \item Carrito persistente en BD para sesión multi-dispositivo.
  \item API REST y microservicios para pagos y notificaciones.
  \item Tests de integración y CI/CD en GitHub Actions.
\end{itemize}

\chapter{Apéndice: fragmentos de código relevantes}

\section{Ejemplo: Transacción de checkout (extracto)}
\begin{lstlisting}[language=PHP]
DB::transaction(function() use ($cart, $user, $validated) {
    foreach ($cart as $item) {
        $product = Product::where('id', $item['id'])
            ->lockForUpdate()->first();
        if ($product->stock < $item['quantity']) {
            throw new Exception('Stock insuficiente');
        }
    }
    $order = Order::create([...]);
    // crear order lines, decrementar stock, crear payment
});
\end{lstlisting}

\chapter{Conclusión}

Esta documentación proporciona una base completa y profesional del proyecto TAD E-commerce Magenta. Si quieres, puedo:
\begin{itemize}
  \item Añadir diagramas simplificados (tablas) para arquitectura y casos de uso.
  \item Generar un PDF (compilar) desde aquí — avisando que la compilación puede fallar si faltan paquetes del sistema.
  \item Eliminar físicamente ficheros de imagen no deseados del repositorio.
\end{itemize}

\vspace{1cm}
Documento preparado por Grupo Magenta — Mayo 2026

\end{document}
